% Source-derived chapter 05: The main argument
% Original source: birational-section-conjecture-strong-birationality
\chapter{The main argument}
\label{birational-section-conjecture-strong-birationality-chapter-05}
\source{birational-section-conjecture-strong-birationality}

\begin{remark}[Source section]
\label{birational-section-conjecture-strong-birationality-chapter-05-source}
\source{birational-section-conjecture-strong-birationality}
\source{birational-section-conjecture-strong-birationality}
This chapter reproduces the corresponding section of the retrieved
LaTeX source, including its definitions, statements, examples, and
proofs. It has not yet been translated into Lean.
\notready
\end{remark}

% The source section title was: The main argument
This section consists of a unique statement. The main theorems will follow easily from it.

\begin{proposition}
\source{birational-section-conjecture-strong-birationality}
\notready\label{simple}
\source{birational-section-conjecture-strong-birationality}
	Let $k$ be a number field, $\nu$ a finite place, write $U=\A^{1}\setminus\{1,2\}$. Let $s\in\Pi_{U/k}(k)$ be a b.l. Galois section and $x_{\nu}\in\P^{1}(k_{\nu})$ the unique associated local point. Let $\delta:\spec k(t)\to\A^{1}$ be the ``diagonal''	 point, i.e. the generic one, and assume that $s_{k(t)}$ lifts to a section of $\spec k(\A^{1})\otimes_{k} k(t)\setminus\{\delta\}$ (e.g. if $s$ is $t$-b.l.). Then $x_{\nu}$ is $k$-rational.
	\begin{proof}
		We may assume that $x_{\nu}\in U$ since otherwise it is clearly $k$-rational. Let $r_{0}$ be the lifting of $s_{k(t)}$ to $\spec k(\A^{1})\otimes_{k} k(t)\setminus\{\delta\}$, and $r$ the image of $r_{0}$ in $U_{k(t)}\setminus \{\delta\}$. We divide the proof in three steps.
		
		\textbf{Step 1. Specializations of $r$.} We are going to prove that, for every $k$-rational point $c\in U(k)$ different from $x_{\nu}$, the specializing loop of $r$ is constant: the idea is to prove that for any specialization $r_{c}$ of $r$ the base change of $r_{c}$ to $k_{\nu}$ is the geometric section associated with $x_{\nu}\in U\setminus\{c\}(k_{\nu})$ regardless of the choice of $r_{c}$, then we apply \autoref{uniqueloop1}.
		
		Write $\Delta\s U\times U$ for the diagonal, consider the family of curves $U\times U\setminus \Delta\to U$ and fix a $k$-rational point $c\in U(k)$ with $c\neq x_{\nu}$, we have $(U\times U\setminus \Delta)_{c}=U\setminus\{c\}$. Consider the scheme $W=\spec\O_{c}\otimes_{k}\O_{c}\setminus\Delta$ over $\spec \O_{c}$, it contains $\spec k(\A^{1})\otimes_{k} k(t)\setminus\{\delta\}$. The morphism $W\to\spec\O_{c}$ is a projective limit of families of curves over $\O_{c}$ (it can be obtained from $U_{\O_{c}}\to\spec \O_{c}$ by removing divisors étale over $\O_{c}$), its generic fiber is $\spec \O_{c}\otimes_{k} k(t)\setminus\{\delta\}$ while the special fiber is $\spec\O_{c}\setminus\{c\}=\spec k(t)$. 
		
		Choose $r_{c}\in\Pi_{U\setminus\{c\}/k}(k)$ any specialization of $r$. Let $r_{1}$ be the image of $r_{0}$ in $\spec \O_{c}\otimes_{k} k(t)\setminus\{\delta\}$. The specializing loop $S^{1}_{c}\to\Pi_{U\setminus\{c\}/k}$ of $r$ factorizes through the specializing loop $S^{1}_{c}\to\Pi_{k(t)/k}$ of $r_{1}$, hence $r_{c}$ is b.l. Furthermore, the image of $r_{c}$ in $U$ is a specialization of $s_{k(t)}$, i.e. it is $s$.
		
		Since $x_{\nu}\in U\setminus\{c\}$, $r_{c}$ is b.l. and maps to $s$, we have that the base change $r_{c,k_{\nu}}$ is \emph{the} geometric section associated to $x_{\nu}$, regardless of the choice of $r_{c}$. By \autoref{uniqueloop1}, the specializing loop of $r$ at $c$ is constant.
	
		%%%%%
	
		\textbf{Step 2. Change of coordinates.} The map $y\mapsto t-y$ defines an automorphism $\phi:\A^{1}_{k(t)}\to\A^{1}_{k(t)}$ with $\phi(\delta)=0$. Write $V=\A^{1}_{k(t)}\setminus\{\phi(1),\phi(2)\}$, we have that $\phi$ restricts to an isomorphism $\phi:U_{k(t)}\to V$. We thus get a Galois section $\phi(s_{k(t)})\in\Pi_{V}(k(t))$ with a lifting $\phi(r)\in\Pi_{V\setminus \{0\}}(k(t))$. Since $V\setminus\{0\}\s\G_{m}$, then $\phi(r)$ induces a section
		\[z\in\Pi_{\G_{m}/k(t)}(k(t))=\cB\hz(1)(k(t))=\projlim_{n} k(t)^*/k(t)^{*n}=\hat{k(t)^*},\]
		\[z=\lambda\cdot\prod_{c}q_{c}^{e_{c}}\]
		where $\lambda\in\hat{k^*}$, $c$ varies among closed points of $\A^{1}$, $q_{c}\in k[t]$ is the monic, irreducible polynomial associated with $c$ and $e_{c}\in\hz$.\footnote{Notice that, while there is an embedding $\hat{k(t)^{*}}\subset \hat{k^{*}}\cdot\prod_{c}q_{c}^{\hz}$, this is not a bijection, since for every $n$ all but finitely many exponents of an element of $\hat{k(t)^{*}}$ are multiples of $n$.} For every $k$-rational point $c\neq x_{\nu},1,2$, the fact that the specializing loop of $r$ at $c$ is constant implies that the same holds for $\phi(r)$ and hence for $z$, hence $e_{c}=0$ by \autoref{gmloop}. Since we can repeat step 1 after base changing to any finite extension of $k$, we get that $e_{c}=0$ for every \emph{closed} point $c\neq x_{\nu},1,2$.
		
		If $x_{\nu}$ is algebraic over $k$, let $q$ be its minimal polynomial and $e$ the exponent of $q$ as a factor of $z$, otherwise $q=1$ and $e=0$. We may thus write
		\[z=\lambda\cdot q^{e}\cdot(t-1)^{e_{1}}\cdot(t-2)^{e_{2}}.\]
		We are going to prove that $q$ has degree $1$ and thus $x_{\nu}$ is $k$-rational.
		
		%%%%%%
		
		\textbf{Step 3. Specializations of $z$.} Fix $K$ a finite extension of $k$ which splits $q$ completely (if $q=1$ choose $K=k$) and let $\eta$ be an extension of $\nu$. In the rest of the proof, we are going to consider many $K$-rational points $c\in K$ (or $c_{i}\in K$). We will always tacitly assume that $c\neq 1,2$ and $q(c)\neq 0$.
		
		If we repeat steps 1 and 2 after base changing to $K$, for every $c$ the base change to $K_{\eta}$ of the unique specialization of $\phi(r)$ at $c$ is $\phi(x_{\eta})$, hence we obtain the equation
		\[\lambda\cdot q(c)^{e}\cdot(c-1)^{e_{1}}\cdot(c-2)^{e_{2}}=c-x_{\eta}\in\hat{K_{\eta}}^{*},\]
		and by applying $\eta:\hat{K_{\eta}}^{*}\to\hz$ we get
		\[\eta(\lambda)+e\eta(q(c))-\eta(c-x_{\eta})+e_{1}\eta(c-1)+e_{2}\eta(c-2)=0\in\hz.\]
		
		Observe that if $p\in K[t]$ is a polynomial, $c\in K$ an element with $p(c)\neq 0$ and $c_{i}\neq c$ is a sequence which tends to $c$ in the $\eta$-adic topology, then $\eta(p(c_{i}))=\eta(p(c))$ is constant for $i>>0$ large enough while $\eta(c_{i}-c)$ is \emph{not} constant.
		
		Choose a sequence of $K$-rational points $c_{i}$ which tends to $1$. Since $1$ is not a root of the polynomials $q$, $t-2$, $t-x_{\eta}$ we see that for $i>>0$ all terms except $e_{1}\eta(c_{i}-1)$ in the equation above are constant. It follows that $e_{1}\eta(c_{i}-1)$ is constant, too. Since $c_{i}\xar{i} 1$ and $c_{i}\neq 1$, then $\eta(c_{i}-1)$ is not constant, this implies that $e_{1}=0$. With the same argument we see that $e_{2}=0$, hence
		\[\eta(\lambda)+e\eta(q(c))-\eta(c-x_{\eta})=0\in\hz.\]
		
		If $x_{\nu}$ is transcendental over $k$ and thus $q=1$, $e=0$, then $\eta(c-x_{\eta})=\eta(\lambda)$ does not depend on $c$, which is absurd, hence $x_{\nu}$ is algebraic over $k$. Since $K$ splits $q$, we may write $q(t)=(t-x_{\eta})\cdot\prod_{j}(t-y_{j})$,
		\[\eta(\lambda)+(e-1)\eta(c-x_{\eta})+e\sum_{j}\eta(c-y_{j})=0\in\hz.\]
		Since we are in characteristic $0$ and $q$ is irreducible over $k$, then $x_{\nu}\neq y_{j}$ for every $j$ and $y_{j}\neq y_{j'}$ for $j\neq j'$. Using a sequence $c_{i}\xar{i}x_{\eta}$ and the same argument as above we see that $e=1$ and hence 
		\[\eta(\lambda)+\sum_{j}\eta(c-y_{j})=0\in\hz.\]
			
		If by contradiction $x_{\nu}$ is not $k$-rational and thus $\deg q\ge 2$, using a sequence $c_{i}\xar{i}y_{j}$ we see that $\eta(\lambda)+\sum_{j}\eta(c-y_{j})=0$ is not constant for $i>>0$, which is absurd.
	\end{proof}
\end{proposition}

%\section{An argument of A. Tamagawa}




%The following is a variation of a famous argument by A. Tamagawa \cite{tam97} \cite[Lemma 53]{sti13} that we are going to use repeatedly. We claim no originality, but we could not find a suitable reference.
	
%\begin{lemma}\label{tama}

%	\begin{proof}
		%~ We do this for quasi-$t$-birational sections, the proof for $t$-birational sections is analogous. We may assume that there is a dominant morphism $f:X\to Z$ with $Z$ non-parabolic such that $f(s)$ is $t$-birational. 
		
		%~ Using \autoref{morb}, we may replace $X$ with a finite étale cover of an open subset and assume that $X$ has genus $\ge 2$.
%		Since $s$ is not geometric nor cuspidal, by the interpretation of cuspidal sections given in \cite[\S 8]{anab} there exists a proper, hyperbolic orbicurve $\mathcal{X}$ with coarse moduli space $\bar{X}$ and an open embedding $j:X\subset\mathcal{X}$ such that $j(s)\in\Pi_{\mathcal{X}/k}$ is not geometric. By \cite[Proposition 7.1]{anab} $\mathcal{X}_{\bar{k}}$ has a finite étale covering by a curve, this implies that $\mathcal{X}_{\bar{k}}\to\Pi_{\mathcal{X}_{\bar{k}}/\bar{k}}$, and hence $\mathcal{X}\to\Pi_{X/k}$, is faithful. Because of this, there exists a finite étale gerbe $\Phi$ with a representable morphism $\mathcal{X}\to\Phi$ with geometrically connected fibers. By the universal property of the étale fundamental gerbe we have a factorization $\mathcal{X}\to\Pi_{\mathcal{X}/k}\to\Phi$, consider the image $\spec k\to\Phi$ of $j(s)$.
		
%		Write $\bar{Y}_{0}=\mathcal{X}\times_{\Phi}\spec k$ and $Y_{0}=X\times_{\Phi}\spec k$, since $\mathcal{X}\to\Phi$ is representable then these are curves which are finite étale coverings of $X$, $\mathcal{X}$ and there is an open embedding $j_{0}:Y_{0}\subset\bar{Y}_{0}$. Since $\mc{X}$ is hyperbolic, $\bar{Y}_{0}$ is hyperbolic too, hence it has genus $\ge 2$ since it is a proper curve. We have an identification $\Pi_{Y_{0}/k}=\Pi_{X/k}\times_{\Phi}\spec k$ and similarly for $\bar{Y}_{0}$, hence $s$, $j(s)$ lift to Galois sections $s_{0}$, $j_{0}(s_{0})$ of $Y_{0}$, $\bar{Y}_{0}$. The fact that $j(s)$ is not geometric implies that $j_{0}(s_{0})$ is not geometric. Since $\bar{Y}_{0}(k)$ is finite, by \cite[Lemma 53]{sti13} we may pass to a finite étale neighbourhood $\bar{Y}\to\bar{Y}_{0}$ with $\bar{Y}(k)=\emptyset$.
%	\end{proof}
%\end{lemma}
